\begin{abstract}
Process supervision assumes that the first incorrect step of a reasoning trace is where it fails,
and verifier-guided search backtracks to that step. We test this on ProcessBench's expert
first-error labels.

Observationally, early-warning detectors mostly read problem difficulty: question and generator
alone reach AUROC \aucQG, and a quarter of the reasoning adds \dIncA\ (n.s.). \nonFatalPct\% of
first errors are not fatal.

We then measure commitment directly. We resample four open models from two families at every
prefix of their own traces (\TotalRollouts\ scored continuations, up to 64 per step in
independent batches) and show that the resulting value curves and decisive steps are reliable.
\QwenSmallAllCapPct--\QwenMathCapPct\% of failures are decided before the first step. When the
outcome is still open, the labelled error is where success collapses in
\QwenTwoAllDec--\QwenSmallAllDec\ of traces, against a null of about 0.18. This was confirmed in a
preregistered test on Llama-3.1-8B.

For three Qwen models, resampling from just before the labelled error succeeds less often than
restarting, and the effect replicates on independent samples (\PooledNewBtRestErr). It does not
generalize to Llama, where the preregistered test fails (\LlamaBtRestErr).

At a matched end-to-end budget that counts every token, no practical policy beats restart
reliably. This includes backtracking guided by a validated process reward model (ProcessBench F1
\prmFOne), which is worse than restart for Llama (\LlamaPolPrmDropFourDiff). A short
supermartingale argument explains why restart is hard to beat.

We release all rollouts, the preregistration and the code behind every number.
\end{abstract}
